%% ══════════════════════════════════════════════════════════════
%%  APPENDIX B — PROMPT BANK SPECIFICATION
%% ══════════════════════════════════════════════════════════════
\chapter{Prompt Bank Specification}
\label{app:prompt-bank}

This appendix documents the structure and content of \autoscholar's prompt bank, which drives autonomous paper generation with 1,000+ deterministically generated research prompts.

\section{Construction Algorithm}
\label{sec:prompts:algorithm}

The prompt bank is built at import time by the function \func{buildPromptBank} in \file{server/services/prompt-bank/index.ts}.
The algorithm interleaves three dimensions:

\begin{algorithm}[H]
\caption{Prompt Bank Construction}
\label{alg:prompt-bank}
\SetAlgoLined
\KwIn{Domains $\mathcal{D}$ (24), Methods $\mathcal{M}$ (24), Topics $\mathcal{T}$ (14 per domain), Compatibility $\mathcal{C}$}
\KwOut{Prompt bank $\mathcal{P}$ with $\leq 1500$ entries}

$\mathcal{P} \gets []$\;
$v \gets 0$\; \tcp{variation index}
\While{$|\mathcal{P}| < 1500$ \textbf{and} $v < 24$}{
  \ForEach{topic index $t \in [0, 13]$}{
    \ForEach{domain $d \in \mathcal{D}$}{
      $\text{methods} \gets \mathcal{C}[d]$\;
      $m \gets \text{methods}[v \bmod |\text{methods}|]$\;
      $\text{topic} \gets \mathcal{T}[d][t]$\;
      $\text{prompt} \gets \text{formatPrompt}(d, m, \text{topic})$\;
      $\mathcal{P}.\text{push}(\{d, m, \text{topic}, \text{prompt}\})$\;
    }
  }
  $v \gets v + 1$\;
}
\Return{$\mathcal{P}$}\;
\end{algorithm}

The iteration order (variation $\to$ topic $\to$ domain) ensures that consecutive prompts always switch domains, providing maximum topic diversity for sequential selection.

\section{Research Domains}
\label{sec:prompts:domains}

\begin{longtable}{rl}
\caption{24 research domains.} \label{tab:domains} \\
\toprule
\textbf{\#} & \textbf{Domain} \\
\midrule
\endfirsthead
\toprule
\textbf{\#} & \textbf{Domain} \\
\midrule
\endhead
1 & Market Microstructure \\
2 & Volatility Modeling \\
3 & Asset Allocation \& Portfolio Management \\
4 & Machine Learning in Finance \\
5 & Macro-Finance \\
6 & Cryptocurrency \& DeFi \\
7 & ESG \& Sustainable Finance \\
8 & Market Anomalies \& Factor Investing \\
9 & Financial Econometrics \\
10 & Systemic Risk \& Financial Stability \\
11 & Behavioral Finance \\
12 & Derivatives \& Options \\
13 & Credit Risk \& Fixed Income \\
14 & Monetary Policy \& Central Banking \\
15 & Banking \& Financial Intermediation \\
16 & Commodities \& Energy Finance \\
17 & Emerging Markets \\
18 & Fintech \& Digital Finance \\
19 & Statistical Inference \\
20 & Machine Learning Prediction \\
21 & Deep Learning \& Neural Networks \\
22 & Applied Statistics \\
23 & Supervised \& Unsupervised Learning \\
24 & AI for Scientific Discovery \\
\bottomrule
\end{longtable}

\section{Methodological Approaches}
\label{sec:prompts:methods}

\begin{longtable}{rl}
\caption{24 methodological approaches.} \label{tab:methods} \\
\toprule
\textbf{\#} & \textbf{Method} \\
\midrule
\endfirsthead
\toprule
\textbf{\#} & \textbf{Method} \\
\midrule
\endhead
1 & Panel Regression with Fixed Effects \\
2 & Generalized Method of Moments (GMM) \\
3 & Event Study Methodology \\
4 & Vector Autoregression (VAR/VECM) \\
5 & DCC-GARCH Multivariate Volatility \\
6 & Difference-in-Differences (DiD) \\
7 & Regression Discontinuity Design (RDD) \\
8 & Instrumental Variables (2SLS/IV) \\
9 & Quantile Regression \\
10 & Stochastic Volatility Models \\
11 & Copula-Based Dependence Modeling \\
12 & Propensity Score Matching \\
13 & Fama-MacBeth Two-Stage Regression \\
14 & Structural Equation Modeling \\
15 & Bayesian Estimation \\
16 & High-Frequency Econometrics \\
17 & Affine Term Structure Models \\
18 & Dynamic Factor Models \\
19 & Bootstrap \& Permutation Testing \\
20 & Random Forest \& Gradient Boosting \\
21 & LSTM \& Transformer Networks \\
22 & Non-Parametric Density Estimation \\
23 & K-Means \& Hierarchical Clustering \\
24 & Reinforcement Learning \\
\bottomrule
\end{longtable}

\section{Domain--Method Compatibility Matrix}
\label{sec:prompts:compat}

Not all domain--method combinations are meaningful.
The compatibility matrix (\file{prompt-bank/domains.ts}) maps each domain to its valid methods.
For example:

\begin{itemize}
  \item \textbf{Market Microstructure:} Panel FE, GMM, Event Study, High-Frequency Econometrics, Quantile Regression, ...
  \item \textbf{Machine Learning in Finance:} Random Forest, LSTM/Transformer, K-Means, Reinforcement Learning, ...
  \item \textbf{Statistical Inference:} Bootstrap, Non-Parametric Density, Bayesian Estimation, Quantile Regression, ...
\end{itemize}

Finance domains are primarily compatible with econometric methods, while statistics/ML/AI domains are compatible with their respective method families.
Cross-pollination is allowed where meaningful (e.g., Bayesian methods for finance, panel regression for ML comparisons).

\section{Topic Examples}
\label{sec:prompts:topics}

Each domain has 14 specific research topics (\file{prompt-bank/topics.ts}, 415 lines).
Examples:

\subsection{Market Microstructure Topics}

\begin{enumerate}
  \item Intraday volatility patterns and market maker behavior
  \item High-frequency trading impact on price discovery
  \item Dark pool trading volume and market quality
  \item Bid-ask spread determinants across market regimes
  \item Flash crash dynamics and circuit breaker effectiveness
  \item Limit order book dynamics and information asymmetry
  \item Market fragmentation effects on execution quality
  \item Algorithmic trading and systemic risk propagation
  \item Price impact of large institutional trades
  \item Cross-market information transmission at the tick level
  \item Maker-taker fee models and liquidity provision incentives
  \item Odd-lot trading and retail investor behavior
  \item Market microstructure during stress events
  \item Co-location advantages and speed competition in equity markets
\end{enumerate}

\subsection{ESG \& Sustainable Finance Topics}

\begin{enumerate}
  \item ESG score divergence across rating agencies
  \item Green bond pricing premiums and issuer characteristics
  \item Carbon risk and asset pricing
  \item ESG momentum and factor investing
  \item Corporate governance and firm value
  \item Climate transition risk in credit markets
  \item Social responsibility and cost of capital
  \item Environmental regulation and competitive advantage
  \item ESG integration in portfolio optimization
  \item Greenwashing detection in corporate disclosures
  \item Biodiversity risk and financial returns
  \item Impact investing performance measurement
  \item Stranded assets and energy transition
  \item ESG controversies and stock price reactions
\end{enumerate}

\section{Prompt Format}
\label{sec:prompts:format}

Each generated prompt follows the structure:

\begin{quote}
\texttt{Write a research paper on [topic] using [method] methodology.
The paper should focus on the [domain] domain and provide rigorous empirical evidence using real financial data.}
\end{quote}

The prompt upgrader (\func{upgradeResearchPrompt}) then expands this into a detailed 400--700 word research specification before it reaches Claude.
